% Q15_gold_answers.tex — GRADER CALIBRATION FIXTURE (correct, paraphrased).
% A drop-in replacement for paper/Q15_answers.tex (spec §4.7, §6.12): every answer is
% correct, but the wording, the seeds and the baseline are deliberately far from the
% reference, so the grader must recognise correctness by substance (the rubric's
% probes), not by matching text. B1 uses a baseline of 0.4 and its own seeds and Claude
% split; the cards word every field their own way. B2 is not graded in this mode.
%
% EXPECTED[B1][hypothesis-precision][4] [budget-parity][4] [variance][4] [contamination][4] [cheat-audit][4] [disclosure][4]: correct. Probe B recomputes every product (720, 480, 80; runs 180, 60) and the totals against the registry (1,200 / 240 and 90 / 20); Probe M at the stated inputs (30, 0.4, paired) prints 0.36564883203538673; Probe W at 3 of 15 prints (0.0705..., 0.4518...); Probe S finds every seed list disjoint; the markers appear once. The report's first line is PREREG Q15: UNLOCKED — every method dimension >= 3. No finding.
% EXPECTED[B3][card-fidelity][4] [B4][card-fidelity][4] [B5][card-fidelity][4]: every field agrees with its pinned range; the Variance and Artifact fields correctly say what each range does not give; the closing lines are correct (B4 takes the other reading of reference-freedom, which line 1675 supports). No finding.
% Expected per-problem grades: B1 = 4, B3 = 4, B4 = 4, B5 = 4 (3 is tolerated only with a declared-gap note and no flaw).
% A calibration run FAILS if any graded problem scores <= 2, if B1 is HELD, or if any
% finding claims a false product, a wrong MDE or Wilson interval, an overlap of seeds or a
% contradiction with a range.
\documentclass[11pt]{article}
\input{preamble}

\title{Q15 --- Practice II-C + Lab 2 \\ \large Written answers}
\author{learner}
\date{}

\begin{document}
\maketitle
\sloppy

% PREREG-BEGIN
\section*{Problem B1 (pre-registration)}\label{ans:B1}
\begin{answer}
\textbf{Claim} Three statements, fixed before any run. (1) Separation: the
verdicts my oracles give agree with the sealed verdicts better than chance
would; I estimate Cohen's $\kappa$ and call the claim refuted if its 95\%
interval reaches 0. (2) Feedback, two-sided: with the same four calls per task,
the share of tasks whose kept specification is faithful is not the same under
S1 as under S0; refuted if the 95\% interval of the paired per-task difference
covers 0. (3) Vericoding, descriptive only: how often V's verified
implementation fails the reference, per language with Wilson intervals; no
comparison with S0 or S1 is made, since V spends eight calls on a different
job.

\textbf{Feedback signals} S1 sees three numbers after each sample: how many of
the shipped implementation's mutants its candidate rejects, whether the
candidate's preconditions can be shown contradictory, and whether its witness
call verifies. All three are computed from the candidate, the task file and
the shipped code; nothing sealed is read, and the implication and
classification oracles are run only after the batch.

\textbf{Metric and estimator} Unit: the task (30 of them). A kept
specification is faithful when the post-batch classification against the
sealed reference says \texttt{FAITHFUL}; anything else, \texttt{UNDECIDED}
included, is unfaithful. For (2): each task's faithful share over its three
repeats, per arm, then the mean of the per-task differences S1 $-$ S0 with a
percentile bootstrap (tasks resampled, 95\%), plus McNemar on each repeat as a
check. For (1): a confusion table of my five verdict classes against the
sealed ones over all kept specifications, $\kappa$ from the table's margins,
bootstrapped over tasks. Verified-but-unfaithful share per arm with a Wilson
interval (for example 3 of 15 gives about 0.07 to 0.45).

\textbf{Arms and policies} S0, S1 and V are three modes of \texttt{render} in
my \texttt{policies.py}; each mode's prompt hash goes in the record once the
file is frozen. Claude runs S1 on 20 tasks, ten per language: 8 local and 2
public in each, drawn with seed 777.

\textbf{Tasks and strata} The local cards listed in
\texttt{bench/local/q/lab2/STRATUM.toml} (ten per language; the file's hash
recorded), five HumanEval-Verus and five HumanEval-Dafny tasks, the public
sets pinned by the commits and validation hashes in
\texttt{bench/manifest.toml}, drawn with seeds 31 and 32; no pilot.

\textbf{Budget arithmetic} S0 + S1: 30 tasks $\times$ 1 attempt $\times$ 4
rounds $\times$ 2 arms $\times$ 3 repeats $= 720$ calls, $30 \times 2 \times 3
= 180$ runs. V: $30 \times 1 \times 8 \times 1 \times 2 = 480$ calls, 60 runs.
Qwen: 1,200 calls, 240 runs. Claude: $20 \times 1 \times 4 \times 1 \times 1 =
80$ calls and up to 10 sentinel calls, 90 calls, 20 runs. This is exactly the
registry's allowance. GPU, from the step-3d throughput in VERSIONS.md (thinking
off, 768-token cap: a median 3.7 s per call at a short prompt, about 0.3 s more
per extra 1k prompt tokens): S0 about 0.37 h; S1 and V at about 8k-token
prompts (about 6 s a call), about 0.6 h and 0.8 h; about 1.8 h in all, inside
3, and only a floor for S1 and V, whose prompts keep growing. Thinking on
would cost about 30 s a call, so it would need far fewer calls.

\textbf{Generation settings} The gate's pinned values for every Qwen call,
identical across arms; seeds S0 501--860, S1 1501--1860, V 2501--2980; Claude
with the pinned argv.

\textbf{Cost per arm} Measured tokens in and out, GPU-seconds and wall clock
per arm against the 3 GPU-hour ceiling, and the share of attempts truncated
per arm as a sensitivity line. Pre-registered fallback: if the measured
per-arm totals would pass the ceiling, S0 and S1 drop to two repeats each.

\textbf{MDE} The comparison pairs 30 tasks;
\texttt{stats.mde\_two\_proportions(30, 0.4, paired=True)} $= 0.366$ at a
baseline of 0.4, power 0.8, $\alpha = 0.05$: only a difference of roughly 37
points is detectable, so a null result is weak evidence.

\textbf{Seeds and repeats} Qwen: the per-request seeds above. Claude:
repeated runs, not seeds.

\textbf{Oracles' version} The sha256 of \texttt{harness/q/fvh\_q/oracles.py}
(\texttt{sha256sum}) at the starting commit, written into each run record.

\textbf{Cheat surface and gate} Every kept specification and every V
implementation passes my Q09 gate core in the verification container (no
holes, shipped code untouched for S0 and S1, frozen statement for V);
classification only after the batch in the grader container; whatever only
that container can decide is reported \texttt{UNKNOWN}.

\textbf{Exclusions} An \texttt{infra-failed} attempt is re-issued with the
same id, prompt hash and seed up to three times, then excluded and counted; a
\texttt{truncated:<cause>} attempt counts as a failure, with a sensitivity
line that drops it; an oracle run that is killed or refused gives
\texttt{unknown}, which is neither a kill nor a pass.

\textbf{Stopping rule} A cap or rate window defers the batch between attempts;
I look at no outcome until the last planned id is scored; no arm is stopped
because of its results.

\textbf{Prediction intervals} (80\%) $\kappa$ from 0.2 to 0.6; S1 $-$ S0
from 0 to 0.3; V's failure rate 0.25--0.65 in Dafny and 0.35--0.75 in Verus.

\textbf{Contamination} The local cards were written on 2026-10-11, after both
models were released. Card contamination audit: a novelty search of each
local card's name and statement, and perturbation probes (four cards with
renamed identifiers, run once per arm, read descriptively). The public sets
are flagged as possibly in training data.

\textbf{Disclosure} This is a small-scale replication with a substituted
model; Claude grades this plan and is also one of the subjects; the oracles are
mine; one machine and one solver pin; each run records model ids and digests,
argv, prompt hashes, the oracles' hash and the sandbox manifest.
\textbf{Batch plan} Prospective file \texttt{research/q15-main.toml}.
This method fixture does not open a batch. Once the policies are final,
record the exact file digest on its own \texttt{plan-sha256:} line and
re-grade the changed pre-registration before executing it. Separately
executed pilot or subset plans each need their own binding and gradebook
row; the registered method stays fixed.
\end{answer}
% PREREG-END

\section*{Problem B2 (lab report)}\label{ans:B2}
\begin{answer}
\textbf{Reproduce:} \verb+<not yet: the runs need the batch runner and the grader container>+
\end{answer}

% ---------------------------------------------------------------------------
\section*{Problem B3 (paper card)}\label{ans:B3}
\begin{answer}
\textbf{Source:} cheatsheet entry 04.6, lines 1618--1649 (the pinned range).

\textbf{Task shape:} Small open models (Qwen-2.5, up to 14B) are trained, SFT
then GRPO, to output a Dafny program with its contract written in, without
reasoning traces.

\textbf{Oracle:} The program must verify; for the subset reward Dafny must also
prove that the generated precondition is no stronger and the generated
postcondition no weaker than the ground truth's.

\textbf{What is trusted:} Dafny and a ground truth produced by Claude 3.5
Sonnet and checked only by the verifier; ``superior'' means superior with
respect to the code.

\textbf{Cheat surface:} Rewarding verification alone produced tautological
specs; the model's full-program output could alter the code.

\textbf{Metric:} Pass@1 verification and SSR: at 14B, RL on the subset reward
gives 85.9\% and 55.3\%, RL on verification alone 92.6\% and 37.3\%.

\textbf{Budget:} Groups of $G = 4$ rollouts per prompt during RL; 128 rollouts
for novel postconditions.

\textbf{Variance:} Nothing on seeds or spread; entropy-bonus runs collapse
after roughly 100 steps.

\textbf{Artifact:} The range lists the data (about 20k programs, DafnyComp) but
no code release.

\textbf{Closest oracle, and whether S1 could receive it:} My
\texttt{implication}, used one way round: it needs the reference, so it is a
grader-side signal, not one S1 may see.
\end{answer}

% ---------------------------------------------------------------------------
\section*{Problem B4 (paper card)}\label{ans:B4}
\begin{answer}
\textbf{Source:} cheatsheet entry 04.7, lines 1653--1677 (the pinned range).

\textbf{Task shape:} Write the contract and invariants of a fixed,
deterministic Dafny method; Qwen2.5 models up to 7B, GRPO with groups of 8,
starting from ReForm's SFT checkpoints.

\textbf{Oracle:} Verification for soundness, and completeness from spectests:
run the real program, have an LLM corrupt the output, and require the
specification to reject the corrupted pair.

\textbf{What is trusted:} Dafny, LLM-generated spectests and executable
SpecCheck predicates.

\textbf{Cheat surface:} Trivial specs, \texttt{assume} and body edits get zero
reward; dropping the verify gate lets the model game the tests.

\textbf{Metric:} On DafnyComp-Spec at 7B, completeness pass@1 rises to 24.30
for SpecRL from 18.15 for SFT, with verifiable pass@1 at 10.40.

\textbf{Budget:} About 58k spectests built offline for roughly \$250.

\textbf{Variance:} The range says no variance or seeds are reported.

\textbf{Artifact:} No code release named.

\textbf{Closest oracle, and whether S1 could receive it:} Closest to
\texttt{kill\_rate}. As built it is not reference-free: the suites come from
the reference specification's weaknesses (line 1675), so S1 could not receive
it unchanged.
\end{answer}

% ---------------------------------------------------------------------------
\section*{Problem B5 (paper card)}\label{ans:B5}
\begin{answer}
\textbf{Source:} cheatsheet entry 03.8, lines 1257--1282 (the pinned range).

\textbf{Task shape:} From a docstring to a runnable postcondition
(\texttt{assert}); GPT-3.5, GPT-4 and StarChat, ten samples per prompt variant.

\textbf{Oracle:} Hidden tests run against a reference implementation for
correctness; GPT-3.5-sampled buggy programs for completeness.

\textbf{What is trusted:} The reference implementation, the test inputs and
the sampled mutants.

\textbf{Cheat surface:} A postcondition that only checks types kills few
mutants (0.27); nothing is proved.

\textbf{Metric:} GPT-4 accept@1 0.77 and accept@10 0.96 on EvalPlus; average
bug-completeness 0.85; 64 Defects4J bugs caught.

\textbf{Budget:} Ten samples per prompt.

\textbf{Variance:} No spread is reported beyond one significance test.

\textbf{Artifact:} The range names no artifact of its own.

\textbf{Closest oracle, and whether S1 could receive it:} \texttt{kill\_rate};
it needs a reference program, not a reference spec, and Lab 2 ships the
program, so S1 may see it.
\end{answer}

\end{document}
